%===================================================================
% Chapter 2: Background
%===================================================================
\chapter{Background}
\label{chap:background}

This chapter provides the shared foundation for the rest of the
dissertation: the query rewriting problem, SQL pipelines and the
extension of rewriting to them, the rule-based paradigm that has
dominated automatic rewriting, the techniques available for
establishing that two queries are equivalent, and the two tools this
dissertation builds on, program synthesis and large language models.
Readers familiar with these topics can skim this chapter and refer
back to it as needed.

\section{The Query Rewriting Problem}
\label{sec:bg:problem}

Throughout this dissertation, a query $q$ is a SQL \texttt{SELECT}
statement over a database whose schema and integrity constraints are
known. Two queries $q$ and $q'$ are \emph{semantically equivalent},
written $q \equiv q'$, if they produce identical results on every
database instance that satisfies those constraints. Identical is meant
under SQL's bag semantics: the same rows with the same multiplicities
(and, for queries whose results are explicitly ordered, in the same
order). The constraints are part of the definition for a practical
reason: many valuable rewrites are correct only because of them. A
rewrite that drops a duplicate-eliminating step, for example, may be
justified by a key constraint and incorrect without it.

The \emph{query rewriting problem} is then: given $q$, find a query
$q'$ such that $q' \equiv q$ and $q'$ executes at lower cost, where
cost means wall-clock latency and, in cloud settings, the money
charged for the computation.

To see where rewriting fits, consider how a database processes a
query. The optimizer first translates the SQL statement into an
\emph{execution plan}, a tree of physical operators such as scans,
joins, and aggregations. Among the many plans that could implement
the given query, it selects the one its internal cost estimates deem
cheapest, and the engine then executes that plan. Most systems
display the selected plan through commands such as \texttt{EXPLAIN},
and the plan shows how the work is distributed across operators;
Chapter~4 relies on this visibility, inspecting a query's plan to
identify its performance bottleneck, the operators where most of the
time is spent.

Optimizers are not limited to plan selection. Most also restructure
the query itself to a degree, applying transformations such as
predicate pushdown and subquery unnesting before and during plan
search. But every transformation an engine applies, at the plan
level or the SQL level, comes from a fixed repertoire that its
developers anticipated, implemented, and shipped. When two
equivalent queries differ in a way that lies outside this
repertoire, the engine optimizes them to very different plans, and
such pairs are common in practice: their running times on the same
engine can differ by orders of magnitude, and every speedup reported
in Chapters~3 through~5 comes, by construction, from a
transformation the underlying engine did not perform on its own.
Query rewriting addresses this gap from outside the engine: it
transforms the SQL before the optimizer sees it, and its output is
ordinary SQL handed to the optimizer like any other query, so
rewriting complements the optimizer rather than replacing it, and
its benefits carry to any engine. One might object that any missing
transformation could simply be added to the engine as one more
rule. That is precisely the rule-based paradigm, and
Section~\ref{sec:bg:rules} examines both its successes and the
structural limits that motivate this dissertation.

Solving the rewriting problem also requires judging cost, which is
itself nontrivial. A query's cost can be judged in three ways: by
executing the query and measuring, which is accurate but expensive;
by asking the engine for an estimate, since the optimizer attaches
its internal cost estimates to the plan it selects, and cloud
warehouses likewise estimate the resources a query will consume
before running it; or by building a dedicated \emph{cost model} that
predicts cost from the query's structure and statistics about the
data, possibly learned from past executions, as the learned cost
model in Chapter~5 does. All three appear in this dissertation.

Whichever technique performs it, rewriting has a price of its own:
search time, verification time, or calls to a commercial language
model. The investment pays off for queries that are executed
repeatedly, where a one-time improvement is harvested on every
subsequent run. Such recurring, resource-intensive queries dominate
production analytics
workloads~\cite{jindal2018selecting,dbtcasestudy1,dbtcasestudy2}, and
they are the setting this dissertation assumes throughout.

\section{SQL Pipelines}
\label{sec:bg:pipelines}

Recurrence at organizational scale takes the form of the \emph{SQL
pipeline}. In tools such as dbt~\cite{dbtWhatExactly}, an
organization's transformation logic is decomposed into SQL
transformations (called \emph{models} in dbt), each defined by a
single \texttt{SELECT} statement, and each transformation declares
which source tables and which other transformations it reads from.
These declarations make the dependencies explicit, so the pipeline
forms a directed acyclic graph: nodes are SQL transformations, and an
edge records that one transformation consumes another's output. An
orchestrator re-executes the pipeline on a schedule, materializing
each node's result as a table or view that downstream nodes,
dashboards, and applications then read. The scale of these pipelines
in production is substantial. Statistics shared by dbt indicate that
five percent of dbt projects, more than a thousand of them, contain
over five thousand models, and Siemens has reported operating more
than 85{,}000 models across 800 dbt projects, with 2{,}500 jobs
executing daily~\cite{dengsoe2023large}. Pipelines of this size are
not written once; they evolve for years with the organizations that
maintain them.

Pipelines enlarge the rewriting problem. The object being rewritten
is now the pipeline itself: individual transformations may be
rewritten as before, but logic may also be moved across nodes, for
example by inlining part of a producer into its consumers or pulling
a predicate from a consumer into a producer; nodes may be split,
merged, added, or removed; and decisions about what to materialize
and how often to refresh become part of the optimization. What must
be preserved also changes. Correctness no longer means preserving
each node's output; it means preserving the pipeline's \emph{final
outputs}, the results that external consumers actually read.
Internal nodes may change or disappear so long as the outputs at the
boundary are unchanged.

In principle, pipeline rewriting reduces to query rewriting. Because
every transformation is itself a \texttt{SELECT} statement, any
final output of a pipeline can be unfolded into a single
self-contained query: substitute each referenced transformation with
its defining statement, and repeat until only source tables remain.
In this sense a pipeline is, semantically, a collection of large
queries, one per final output, and pipeline equivalence is exactly
query equivalence applied to each of them. This reduction supplies
the correctness criterion stated above. As a representation for
optimization, however, the unfolded query discards precisely the
structure that matters. First, it hides sharing: a transformation
read by several outputs reappears as duplicated logic inside each
unfolded query, so the central pipeline-level question, namely what
to compute once, persist, and reuse across consumers, cannot even be
posed within any single query. Second, it collapses time: a pipeline
persists intermediate results across runs and refreshes different
nodes at different frequencies, while a query describes one
execution and has no notion of persistent intermediate state or of
parts running on different schedules. Third, it fails at scale:
unfolding a real project produces deeply nested SQL with tens or
hundreds of layers, far beyond what today's optimizers, rewriters,
or equivalence checkers can process, and since the engine optimizes
each statement in isolation, materialization boundaries are walls
its optimizer cannot see across in any case. The dependency graph is
therefore not incidental packaging around one large query; it is the
representation that keeps sharing, persistence, and refresh
explicit, and those are exactly the dimensions in which
pipeline-level optimization lives. Chapter~5 builds directly on this
structure.

\section{Rule-Based Query Rewriting}
\label{sec:bg:rules}

For four decades, the dominant approach to automatic query rewriting
has been the \emph{rewrite rule}: a syntactic pattern, a condition
under which the transformation is safe, and a replacement. When a
fragment of a query matches the pattern and satisfies the condition,
the fragment is replaced. Classical rules include pushing filtering
predicates closer to the data they filter~\cite{levy1994query},
unnesting nested subqueries into joins~\cite{kim1982optimizing}, and
eliminating redundant joins or duplicate-removal steps that
constraints render unnecessary.

Rule-based rewriting entered database systems with
Starburst~\cite{pirahesh1992extensible}, which made rewriting an
extensible phase of query processing, and it remains the architecture
of modern optimizers. Apache Calcite~\cite{begoli2018apache}, the
open-source optimizer embedded in many data systems, ships with
hundreds of rules. The paradigm's known weaknesses have each drawn
research attention: because writing correct rules demands rare
expertise, WeTune~\cite{wetune} discovers and verifies new rules
automatically; because applying rules in the wrong order can miss
improvements or degrade performance, LearnedRewrite~\cite{LR} learns
which rules to apply and in what order, guided by a cost model.
Calcite and LearnedRewrite serve as rule-based points of comparison
in the chapters that follow.

These efforts strengthen the paradigm without changing its boundary.
However a rule set is obtained, curated or discovered, well ordered
or poorly ordered, a rule-based rewriter can express only the
rewrites its rules encode. Each rule must be anticipated, specified,
and proven before it helps its first query, and a query whose
improvement lies outside the rule set receives nothing. The chapters
that follow treat rule-based systems both as baselines to compare
against and as the boundary to move beyond.

\section{Establishing Query Equivalence}
\label{sec:bg:equivalence}

Every rewriting technique in this dissertation must answer the same
question: is the candidate $q'$ really equivalent to $q$? For
rule-based systems the question is settled once per rule, when the
rule is proven. A system that generates rewrites freely must instead
settle it once per candidate, which makes equivalence checking part
of the rewriting process itself. The question is also fundamentally
hard: deciding the equivalence of relational queries is undecidable
in general, a classical consequence of Trakhtenbrot's
theorem~\cite{trakhtenbrot1950}, and deciding even restricted
fragments is expensive. Practical systems therefore combine several
techniques, trading coverage, cost, and strength of assurance
against one another.

\emph{Testing} executes both queries on concrete database instances
and compares their results. A single mismatching instance is a
\emph{counterexample} that definitively refutes equivalence, and a
counterexample is useful beyond refutation: it can steer a search
away from a mistaken region, or show a language model precisely how
its suggestion fails. Agreement on all tested instances, however,
proves nothing about untested ones, so testing serves as a fast,
incomplete filter that eliminates most incorrect candidates cheaply.

\emph{Formal verification} constructs a proof. Bounded verification
proves two queries equivalent on all database instances up to a size
bound, typically by encoding the check as constraints for an SMT
solver; a proof at bound $k$ leaves open what happens beyond $k$,
but inequivalent query pairs are overwhelmingly likely to differ
already on small instances, an assumption borne out empirically in
Chapter~3. Full verification aims at unconditional proofs, through
algebraic reasoning and automated
provers~\cite{chu2017cosette,chu2018axiomatic,zhou2022spes}. Full
verifiers provide the strongest assurance and are the most
restricted in the SQL they accept, which is why practical systems
layer these techniques, spending cheap tests first and reserving
expensive proofs for the few candidates that remain.

At pipeline scale, formal proof is out of reach: a refactoring may
touch dozens of interdependent transformations, and existing
verifiers support neither the SQL features nor the multi-statement
scope involved. \emph{Data-backed checking} takes the pragmatic
point on the spectrum: execute the original and refactored pipelines
on production-scale data and compare the materialized outputs that
external consumers read. The resulting assurance is relative to the
data tested rather than universal, but it concerns exactly the
outputs the organization depends on, on exactly the data it runs.

The three systems in this dissertation each assemble their assurance
differently. SlabCity (Chapter~3) layers testing beneath bounded and
full verification inside its search loop. GenRewrite (Chapter~4)
uses testing to catch inequivalent suggestions and turns the
resulting counterexamples into corrections. DAGSmith (Chapter~5)
admits a pipeline refactoring only after data-backed checking
confirms its final outputs.

\section{Program Synthesis and Large Language Models}
\label{sec:bg:tools}

Moving beyond rules requires a source of rewrites that no one wrote
down in advance. This dissertation draws on two.

\emph{Program synthesis} is the task of automatically finding a
program, within a defined space, that satisfies a specification.
Query rewriting instantiates it naturally: the space is SQL, and the
specification is the original query itself, since the synthesized
program must be equivalent to it and faster. What makes synthesis
hard is scale: the space of candidate programs grows explosively
with program size, so practical synthesizers combine enumeration
with aggressive pruning, discarding whole families of candidates at
once, and with feedback loops such as counterexample-guided
inductive synthesis, in which each failed candidate's counterexample
sharpens the tests that future candidates must
pass~\cite{solar2008program}. Synthesis has been applied to SQL
chiefly to generate queries from input-output
examples~\cite{wang2017synthesizing}; using it to optimize an
existing query, where the specification is another query rather than
examples, is the step Chapter~3 takes.

\emph{Large language models} are neural networks trained on vast
corpora of text and code to predict the next token, a simple
objective that at sufficient scale yields broad competence in
understanding and generating programs, SQL
included~\cite{zhao2023llm}. They are used through \emph{prompting}:
a task description, optionally with examples, supplied as input
text~\cite{googleprompt}. Their operational costs are measured in
\emph{tokens}, the subword units in which they consume and produce
text, and commercial models charge per token~\cite{openaitokens},
so the number and length of model calls is a first-class engineering
concern. For rewriting, their promise is exactly the experience that
rules cannot capture: trained on decades of public code and
discussion, they recognize how experienced engineers reshape
queries, including patterns no rule set anticipates. Their weakness
is reliability. A model produces fluent, confident output with no
guarantee of correctness, a failure mode commonly called
hallucination, and it has no innate sense of whether a suggested
rewrite is faster. Every use of language models in this dissertation
is therefore paired with the checking machinery of
Section~\ref{sec:bg:equivalence}: models propose candidates, and
independent checking decides which survive.